\section{系统架构与数据通路}

\subsection{多模态感知的层级}

GR00T 架构由三个视觉/语言/力感知的层级组成：

\begin{itemize}
    \item \textbf{视觉编码器}：处理 RGB-D、lidar 扫描，输出空间感知 embedding
    \item \textbf{语言程序器}：将指令解析成动作意图、工具调用序列、对齐 reward
    \item \textbf{动作调度器}：根据动力学限制选择 torque/force/action，实时闭环
\end{itemize}

\begin{knowledgebox}{感知-语言-动作三段论}
这三个层级形成一个感知-意图-执行的闭环：视觉告诉语言模块“场景有什么”，语言模块告诉动作模块“我要做什么”，动作模块反馈结果并根据 reward 更新策略。
\end{knowledgebox}

\subsection{数据流与训练路径}

典型训练路径如下：

\begin{enumerate}
    \item 收集 teleoperation/telemetry 数据（人类操作 + sensor log）
    \item 用 LLM 自动生成任务 prompt（Eureka pipeline）
    \item 在 Isaac Sim 中扩展任务并生成 synthetic replay
    \item 对 replay 施加 offline RL（off-policy data + behavior cloning）
    \item 用 sim-to-real gap alignment（domain randomization + system ID）
\end{enumerate}

\begin{importantbox}{Replay Buffer hygiene}
在 replay buffer 中保留多样性：最近成功 trajectories、少数失败样本、expert demonstration、LLM-generated exploration。顺序 replay + prioritized replay 结合可以防止模型过度拟合常见样本。
\end{importantbox}

\subsection{Sim-to-Real Alignment 细节}

GR00T 的桥接策略采用三重保障：

\begin{enumerate}
    \item \textbf{Domain randomization}：扰动光照、质地、摩擦系数，让模型适应更宽泛的物理场景
    \item \textbf{System identification}：利用少量 real-world trajectory 调整仿真参数，确保动力学相符
    \item \textbf{Validation loop}：在真实 robot 上运行简化任务，将分布外次数记录到 DevOps dashboard
\end{enumerate}

\begin{importantbox}{Sim-to-Real leash}
每个 deployment 都附带一个 ``sim-to-real leash''：只要 Validation loop 没通过，就保持模型在 sandbox 模式，直到追踪到 tolerable drift level 才解除。
\end{importantbox}

\subsection{本章小结}

Sim-to-Real alignment 需要 domain randomization、system identification 和 validation loop 三者形成闭环；leash 与 guard rail 保证仿真-真实之间的渐进过渡。

